\chapter{Operational Triage: OODA Loop and MITRE ATT\&CK Mapping}
\label{ch:mitre}

A machine-learning intrusion detector does not run in isolation. It
feeds into a SOC staffed by human analysts, integrates with a SIEM, and
is ultimately judged on whether it helps the organisation respond to
real incidents. This chapter sketches how the artefacts produced by
this thesis fit into a SOC workflow, organised by Boyd's OODA loop, and
documents the MITRE ATT\&CK mapping that turns model output into
playbook-ready alerts.

\section{The OODA Loop in a SOC Workflow}

Boyd's OODA loop --- Observe, Orient, Decide, Act --- describes the
decision cycle of any adversarial engagement. Whoever cycles faster
wins. In a SOC context, ``Observe'' is alert ingestion, ``Orient'' is
enrichment with context (reputation lookups, asset criticality, MITRE
mapping), ``Decide'' is the analyst's classification of the alert as a
real incident or a false positive, and ``Act'' is containment and
remediation.

In this thesis, the trained binary detector lives at the Observe stage:
the flow extractor ingests packets, the model scores each flow, and an
alert is emitted when the predicted probability exceeds the Youden-J
threshold. Inference latency at batch~=~10{,}000 is 3.2~$\mu$s/flow for
random forest --- well under any operational requirement. The MITRE
mapping described in Section~\ref{sec:mitre-mapping} provides the Orient
enrichment that turns the alert from ``probability 0.97'' into a
structured story: ``FTP-Patator on port 21, T1110.001 Brute Force,
Credential Access, from external IP X with 1{,}400 similar alerts in the
last hour''.

The analyst's Decide step --- escalate or close --- shrinks under this
enrichment because the alert routes directly to a known playbook (T1110
to credential-stuffing playbook; T1498 to DDoS playbook). The detector
itself does not Act; it produces evidence for action, which is taken by
humans or by automation rules tied to playbook outcomes. This separation
is intentional: a detector with a 21\%--24\% adversarial evasion rate
cannot be trusted to make blocking decisions autonomously.

The bottleneck in OODA is the analyst's Decide step (microseconds for
the detector, sub-second for the orchestrator, minutes for the human).
The MITRE mapping is therefore the single largest force multiplier in
this design, because it removes the manual ``which playbook applies''
step.

\section{Mapping Predictions to MITRE ATT\&CK}
\label{sec:mitre-mapping}

MITRE ATT\&CK~\cite{Strom2018} is the industry-standard taxonomy for
adversary behaviour, organised by tactic (the ``why'') and technique
(the ``how''). It is more granular than the Lockheed kill chain (over
200 techniques) and is maintained as an open knowledge base. SOC
playbooks are usually keyed to ATT\&CK technique IDs, so mapping
classifier output to ATT\&CK lets the SOC route alerts to existing
playbooks rather than triaging each prediction in isolation.

\begin{figure}[H]
    \centering
    \includegraphics[width=0.85\textwidth]{tactic_distribution.png}
    \caption{Distribution of CICIDS2017 attack flows across MITRE ATT\&CK
    tactics. The dataset is dominated by Impact (DDoS, DoS Hulk) and
    Credential Access (Patator) tactics; Initial Access techniques such
    as Heartbleed contribute very few flows.}
    \label{fig:tactic-distribution}
\end{figure}

The binary random-forest detector produces high-precision alerts on
T1498 (Network Denial of Service: DDoS and DoS Hulk) and T1046 (Network
Service Discovery: PortScan). These attack types have distinctive flow
signatures --- short, bursty, asymmetric for DoS; many short flows to
many ports for PortScan --- and the detector catches them at low FPR
with high precision.

Multi-class predictions on the day split are unusable for sub-class
attribution because the model cannot output a label it has not seen.
Heartbleed (T1190) has only 11 examples in the dataset; Web Attack
subtypes are confused with each other. The deployed system therefore
uses the binary head to flag, and the multi-class head only when the
predicted class is in-distribution.

\begin{table}[H]
    \centering
    \caption{Predicted class to MITRE ATT\&CK mapping.}
    \label{tab:mitre-mapping}
    \footnotesize
    \begin{tabular}{p{4.0cm}p{1.8cm}p{5.0cm}p{3.5cm}}
        \toprule
        \textbf{Predicted class} & \textbf{ATT\&CK ID} & \textbf{Technique} & \textbf{Phase} \\
        \midrule
        FTP-Patator & T1110.001 & Brute Force: Password Guessing & Credential Access \\
        SSH-Patator & T1110.001 & Brute Force: Password Guessing & Credential Access \\
        Web Attack-Brute Force & T1110 & Brute Force & Credential Access \\
        DDoS & T1498 & Network Denial of Service & Impact \\
        DoS Hulk / GoldenEye & T1498 & Network Denial of Service & Impact \\
        DoS Slowhttptest / slowloris & T1499 & Endpoint Denial of Service & Impact \\
        PortScan / Infiltration & T1046 & Network Service Discovery & Discovery \\
        Bot & T1071.001 & Application Layer Protocol: Web Protocols & Command and Control \\
        Heartbleed / SQLi & T1190 & Exploit Public-Facing Application & Initial Access \\
        Web Attack-XSS & T1059.007 & Command and Scripting: JavaScript & Execution \\
        \bottomrule
    \end{tabular}
\end{table}

The mapping is static for two reasons. First, a learned mapping would
require richer telemetry than this thesis has access to (logs, EDR,
identity signals). Second, the static mapping is auditable and portable:
a successor SOC can read the JSON file, understand exactly why a given
prediction routes to a given playbook, and modify the mapping if their
attack catalogue or playbook taxonomy differs.
